\documentclass{beamer}
\usetheme{default}
\usebeamercolor{whale}
\title{\textbf{TITLE}}
\subtitle{SUBTITLE}
\author{YUVARAJ GEORGE}
\institute[BITS Pilani, Hyderabad Campus]{Department of Mechanical Engineering\\BITS Pilani, Hyderabad Campus}
\date{September 2021}


\usepackage{ragged2e}  % `\justifying` text
\usepackage{booktabs}  % Tables
\usepackage{tabularx}
\usepackage{tikz}      % Diagrams
\usetikzlibrary{calc, shapes, backgrounds}
\usepackage{amsmath, amssymb}
\usepackage{url}       % `\url`s
\usepackage{listings}  % Code listings
\frenchspacing
\usepackage{graphicx}



\defbeamertemplate*{title page}{customized}[1][]
{
  \usebeamerfont{title}\inserttitle\par
  \usebeamerfont{subtitle}\usebeamercolor[fg]{subtitle}\insertsubtitle\par
  \bigskip
  \usebeamerfont{author}\insertauthor\par
  \usebeamerfont{Department}\insertinstitute\par
  \usebeamerfont{date}\insertdate\par
  \usebeamercolor[fg]{titlegraphic}\inserttitlegraphic
}

\setbeamertemplate{footline}
{
  \leavevmode%
  \hbox{%
  \begin{beamercolorbox}[wd=.25\paperwidth,ht=2.25ex,dp=1ex,center]{author in head/foot}%
    \usebeamerfont{author in head/foot}\insertshortinstitute
  \end{beamercolorbox}%
  \begin{beamercolorbox}[ht=2.25ex,dp=1ex,center]{title in head/foot}%
    \usebeamerfont{title in head/foot}\insertshorttitle\hspace*{1.5em}
    \usebeamerfont{title in head/foot}\insertdate\hspace*{1.5em}
    \insertframenumber{} / \inserttotalframenumber\hspace*{0.8ex}
  \end{beamercolorbox}}%
  \vskip0pt%
}
\makeatletter
\setbeamertemplate{navigation symbols}{}

\begin{document}

{\usebackgroundtemplate{%
 \begin{picture}(100,220)
  \includegraphics[width=\paperwidth,height=\paperheight,keepaspectratio]{Slide1.JPG}\end{picture}} 
  \frame{\maketitle}
{\usebackgroundtemplate{%
  \includegraphics[width=\paperwidth,height=\paperheight,trim= 0 12 0 0 cm,clip=true]{Slide2.JPG}} 
\AtBeginSection[]{% Print an outline at the beginning of sections
    \begin{frame}<beamer>
      \frametitle{Outline for Section \thesection}
      \tableofcontents[currentsection]
    \end{frame}}

\section{Introduction}
    \subsection{Blind Text}
    \begin{frame}{Just like it}
      \framesubtitle{Frame subtitle 1}%
      \begin{tikzpicture}[overlay,remember picture]
        \node[anchor=south east,xshift=-30pt,yshift=35pt]
          at (current page.south east) {
            \includegraphics[width=35mm]{1.jpg}
          };
      \end{tikzpicture}%
      'The song I came to sing\\
remains unsung to this day.\\
I have spent my days in stringing\\
and in unstringing my instrument.\\\bigskip

      The time has not come true,\\
the words have not been rightly set;\\
only there is the agony\\
of wishing in my heart…..\\ \bigskip

\hspace{2 em} - Tagore

    \end{frame}

    \begin{frame}[label=lists]{Itemize \& Enumerate}
      \framesubtitle{Lorem ipsum dolor sit amet}
      \begin{columns}[onlytextwidth]
        \column{.5\textwidth}
          \begin{itemize}
            \item Nulla nec lacinia odio. Curabitur urna tellus.
            \begin{itemize}
              \item Fusce id sodales dolor. Sed id metus dui.
              \begin{itemize}
                \item Cupio virtus licet mi vel feugiat.
              \end{itemize}
            \end{itemize}
          \end{itemize}
        \column{.5\textwidth}
          \begin{enumerate}
            \item Donec porta, risus porttitor egestas scelerisque video.
            \begin{enumerate}
              \item Nunc non ante fringilla, manus potentis cario.
              \begin{enumerate}
                \item Pellentesque servus morbi tristique.
              \end{enumerate}
            \end{enumerate}
          \end{enumerate}
      \end{columns}
      \bigskip
      \justifying

      {\uselanguage{czech}Nechť již hříšné saxofony ďáblů
      rozzvučí síň úděsnými tóny waltzu, tanga a quickstepu!}
      {\uselanguage{slovak} Nezvyčajné kŕdle šťastných figliarskych
      ďatľov učia pri kótovanom ústí Váhu mĺkveho koňa Waldemara
      obžierať väč\-šie kusy exkluzívnej kôry.}
      {\uselanguage{english}The quick, brown fox jumps over a lazy
      dog. DJs flock by when MTV ax quiz prog. “Now fax quiz Jack!”}
    \end{frame}

    \subsection{Structuring Elements}
    \begin{frame}[label=simmonshall]{Text blocks}
      \framesubtitle{In plain, example, and \alert{alert} flavour}
      \alert{This text} is highlighted.

      \begin{block}{A plain block}
        This is a plain block containing some \alert{highlighted text}.
      \end{block}
      \begin{exampleblock}{An example block}
        This is an example block containing some \alert{highlighted text}.
      \end{exampleblock}
      \begin{alertblock}{An alert block}
        This is an alert block containing some \alert{highlighted text}.
      \end{alertblock}
      
      can define additional environments in preamble
    \end{frame}

   

    \subsection{Mathematical Expressions}
    \begin{frame}[label=math]{Numerals and Mathematics}
      \framesubtitle{Formulae, equations,etc}
      \begin{equation}
         \frac{1}{2} \rho V^2 S C_{L} = \Vec{L}
      \end{equation}
      
      \begin{equation}
         \Gamma =\oint _{{C}}V\cdot d{\mathbf  {s}}=\oint_{{C}}V\cos \theta 
      \end{equation}
    \end{frame}

    \subsection{Figures and Code Listings}
    \begin{frame}[label=figs1]{Figures}
      \framesubtitle{Tables, graphs, and images}
      \begin{table}
        \begin{tabular}{ccc}
        \toprule
          \textbf{S.NO} & \textbf{parameter} & \textbf{Value}\\
          \toprule
         1 & X & 1\\
         2 & P & 5\\
         3 & n! & 6666\\
          \bottomrule
        \end{tabular}
        \caption{an example showing tables}
      \end{table}
    \end{frame}
    \begin{frame}[label=figs2]{Figures}
      \framesubtitle{Tables, graphs, and images}
      \begin{figure}[b]
        \centering
        % Flipping a coin
        % Author: cis
        \tikzset{
          head/.style = {fill = none, label = center:\textsf{H}},
          tail/.style = {fill = none, label = center:\textsf{T}}}
        \scalebox{0.65}{\begin{tikzpicture}[
            scale = 1.5, transform shape, thick,
            every node/.style = {draw, circle, minimum size = 10mm},
            grow = down,  % alignment of characters
            level 1/.style = {sibling distance=3cm},
            level 2/.style = {sibling distance=4cm},
            level 3/.style = {sibling distance=2cm},
            level distance = 1.25cm
          ]
          \node[shape = rectangle,
            minimum width = 6cm, font = \sffamily] {Coin flipping}
          child { node[shape = circle split, draw, line width = 1pt,
                  minimum size = 10mm, inner sep = 0mm, rotate = 30] (Start)
                  { \rotatebox{-30}{H} \nodepart{lower} \rotatebox{-30}{T}}
           child {   node [head] (A) {}
             child { node [head] (B) {}}
             child { node [tail] (C) {}}
           }
           child {   node [tail] (D) {}
             child { node [head] (E) {}}
             child { node [tail] (F) {}}
           }
          };

          % Filling the root (Start)
          \begin{scope}[on background layer, rotate=30]
            \fill[head] (Start.base) ([xshift = 0mm]Start.east) arc (0:180:5mm)
              -- cycle;
            \fill[tail] (Start.base) ([xshift = 0pt]Start.west) arc (180:360:5mm)
              -- cycle;
          \end{scope}

          % Labels
          \begin{scope}[nodes = {draw = none}]
            \path (Start) -- (A) node [near start, left]  {$0.5$};
            \path (A)     -- (B) node [near start, left]  {$0.5$};
            \path (A)     -- (C) node [near start, right] {$0.5$};
            \path (Start) -- (D) node [near start, right] {$0.5$};
            \path (D)     -- (E) node [near start, left]  {$0.5$};
            \path (D)     -- (F) node [near start, right] {$0.5$};
            \begin{scope}[nodes = {below = 11pt}]
              \node [name = X] at (B) {$0.25$};
              \node            at (C) {$0.25$};
              \node [name = Y] at (E) {$0.25$};
              \node            at (F) {$0.25$};
            \end{scope}
          \end{scope}
        \end{tikzpicture}}
        \caption{Tree of probabilities -- Flipping a coin\footnote[frame]{%
          A derivative of a diagram from \url{texample.net} by cis, CC BY 2.5 licensed}}
      \end{figure}
    \end{frame}

    \defverbatim[colored]\sleepSort{
      \begin{lstlisting}[language=C,tabsize=2]
  #include <stdio.h>
int main() {
   // printf() displays the string inside quotation
   printf("Hello, World!");
   return 0;
}
    \end{lstlisting}}
    \begin{frame}{Code listings}{An example source code in C}
      \sleepSort
    \end{frame}

    \subsection{Citations and Bibliography}
    \begin{frame}[label=citations]{Citations}
      \framesubtitle{\TeX, \LaTeX, and Beamer}
\hspace{1 em}
      \justifying\TeX\ is a programming language for the typesetting
      of documents. It was created by Donald Erwin Knuth in the late
      1970s and it is documented in \emph{The \TeX
      book}~\cite{knuth84}.

      In the early 1980s, Leslie Lamport created the initial version
      of \LaTeX, a high-level language on top of \TeX, which is
      documented in \emph{\LaTeX : A Document Preparation
      System}~\cite{lamport94}. There exists a healthy ecosystem of
      packages that extend the base functionality of \LaTeX;
      \emph{The \LaTeX\ Companion}~\cite{MG94} acts as a guide
      through the ecosystem.  In 2003, Till Tantau created the initial version of Beamer, a
      \LaTeX\ package for the creation of presentations. Beamer is
      documented in the \emph{User's Guide to the Beamer
      Class}~\cite{tantau04}.
    \end{frame}

    \begin{frame}[label=bibliography]{Bibliography}
      \framesubtitle{\TeX, \LaTeX, and Beamer}
      \begin{thebibliography}{9}
        \bibitem{knuth84}
            Donald~E.~Knuth.
            \emph{The \TeX book}.
            Addison-Wesley, 1984.
        \bibitem{lamport94}
            Leslie~Lamport.
            \emph{\LaTeX : A Document Preparation System}.
            Addison-Wesley, 1986.
        \bibitem{MG94}
            M.~Goossens, F.~Mittelbach, and A.~Samarin.
            \emph{The \LaTeX\ Companion}.
            Addison-Wesley, 1994.
        \bibitem{tantau04}
            Till~Tantau.
            \emph{User's Guide to the Beamer Class Version 3.01}.
            Available at \url{http://latex-beamer.sourceforge.net}.
        \bibitem{MS05}
            A.~Mertz and W.~Slough.
            Edited by B.~Beeton and K.~Berry.
            \emph{Beamer by example} In TUGboat,
              Vol. 26, No. 1., pp. 68-73.
      \end{thebibliography}
    \end{frame}

\end{document}

\end{document}
